\FigureLayoutDeclare{img/2014/hubei_science/q5_fig1.png}{13.0cm}{18bp 31bp 140bp 16bp}

\chapter{2014年湖北卷（理）}

\section{单选题}

\begin{problem}
\(\i\) 为虚数单位，则 \(\left(\frac{1 - \i}{1 + \i}\right)^2 =\)（\quad）

\choices
    {\(-1\)}
    {\(1\)}
    {\(-\i\)}
    {\(\i\)}
\end{problem}

\begin{answer}
A
\end{answer}

\begin{solution}
\(1-\i\) 与 \(1+\i\) 互为共轭复数，因此
\[
\frac{1-\i}{1+\i}=\frac{(1-\i)^2}{(1+\i)(1-\i)}=\frac{-2\i}{2}=-\i
\]
所以原式 \(=(-\i)^2=-1\)，故选 A.
\end{solution}

\begin{problem}
若二项式 \(\left(2x + \frac{\alpha}{x}\right)^7\) 的展开式中 \(\frac{1}{x^3}\) 的系数是 \(84\)，则实数 \(\alpha =\)（\quad）

\choices
    {\(2\)}
    {\(\sqrt[3]{4}\)}
    {\(1\)}
    {\(\frac{\sqrt[3]{2}}{4}\)}
\end{problem}

\begin{answer}
C
\end{answer}

\begin{solution}
展开式的通项为 \(T_{k+1}=\mathrm C_7^k(2x)^{7-k}\left(\frac{\alpha}{x}\right)^k\). 因而 \(x\) 的指数为 \(7-2k\). 要得到 \(x^{-3}\)，必须有 \(7-2k=-3\)，故 \(k=5\)，此时 \(x^{-3}\) 的系数为
\[
\mathrm C_7^5\,2^2\alpha^5=84\alpha^5
\]
由题意 \(84\alpha^5=84\)，得 \(\alpha^5=1\). 因为 \(\alpha\) 为实数，所以 \(\alpha=1\)，故选 C.
\end{solution}

\begin{problem}
设 \(U\) 为全集，\(A\)，\(B\) 是集合，则“存在集合 \(C\) 使得 \(A \subseteq C\)，\(B \subseteq \complement_U C\)”是 \(A \cap B = \varnothing\) 的（\quad）

\choices
    {充分而不必要条件}
    {必要而不充分条件}
    {充要条件}
    {既不充分也不必要条件}
\end{problem}

\begin{answer}
C
\end{answer}

\begin{solution}
若存在集合 \(C\) 满足 \(A\subseteq C\)，\(B\subseteq\complement_U C\)，则 \(A\) 与 \(B\) 不可能有公共元素，即 \(A\cap B=\varnothing\). 反之，若 \(A\cap B=\varnothing\)，取 \(C=A\)，则 \(A\subseteq C\)，且 \(B\subseteq\complement_U A=\complement_U C\). 因此两个条件互为充要条件，故选 C.
\end{solution}

\begin{problem}
根据如下样本数据

\begin{center}
\begin{tabular}{c|cccccc}
\hline
\(x\) & \(3\) & \(4\) & \(5\) & \(6\) & \(7\) & \(8\) \\
\hline
\(y\) & \(4.0\) & \(2.5\) & \(-0.5\) & \(0.5\) & \(-2.0\) & \(-3.0\) \\
\hline
\end{tabular}
\end{center}

得到的回归方程为 \(\hat{y} = bx + a\)，则（\quad）

\choices
    {\(a > 0\)，\(b > 0\)}
    {\(a > 0\)，\(b < 0\)}
    {\(a < 0\)，\(b > 0\)}
    {\(a < 0\)，\(b < 0\)}
\end{problem}

\begin{answer}
B
\end{answer}

\begin{solution}
由表中数据得 \(\overline{x}=\frac{3+4+5+6+7+8}{6}=\frac{11}{2}\)，\(\overline{y}=\frac{4.0+2.5-0.5+0.5-2.0-3.0}{6}=\frac14\). 回归系数 \(b\) 的符号与 \(\sum (x_i-\overline{x})(y_i-\overline{y})\) 的符号相同，而
\[
\sum (x_i-\overline{x})(y_i-\overline{y})=\sum (x_i-\overline{x})y_i=-\frac{95}{4}<0
\]
所以 \(b<0\). 进一步 \(a=\overline{y}-b\overline{x}=\frac14-\frac{11}{2}b>0\). 故选 B.
\end{solution}

\begin{problem}
在如图所示的空间直角坐标系 \(O-xyz\) 中，一个四面体的顶点坐标分别是 \((0,0,2)\)，\((2,2,0)\)，\((1,2,1)\)，\((2,2,2)\). 给出编号为\(\circled{1}\)、\(\circled{2}\)、\(\circled{3}\)、\(\circled{4}\)的四个图，则该四面体的正视图和俯视图分别为（\quad）

\bitmapfigure[max width=0.42\linewidth]{img/2014/hubei_science/q5_fig1.png}

\choices
    {\(\circled{1}\)和\(\circled{2}\)}
    {\(\circled{3}\)和\(\circled{1}\)}
    {\(\circled{4}\)和\(\circled{3}\)}
    {\(\circled{4}\)和\(\circled{2}\)}
\end{problem}

\begin{answer}
D
\end{answer}

\begin{solution}
记四个顶点依次为 \(V_1=(0,0,2)\)、\(V_2=(2,2,0)\)、\(V_3=(1,2,1)\)、\(V_4=(2,2,2)\). 正视图是沿 \(x\) 轴方向投影到 \(yOz\) 平面，于是
\(V_1'=(0,2)\)，\(V_2'=(2,0)\)，\(V_3'=(2,1)\)，\(V_4'=(2,2)\).
其中 \(V_2'\)、\(V_3'\)、\(V_4'\) 共线，四面体的外轮廓为三角形 \(V_1'V_2'V_4'\)，棱 \(V_1V_3\) 的投影为从 \(V_1'\) 到 \(V_3'\) 的虚线，故正视图是图\(\circled{4}\).

俯视图是沿 \(z\) 轴方向投影到 \(xOy\) 平面，四个顶点的投影为
\(V_1''=(0,0)\)，\(V_2''=(2,2)\)，\(V_3''=(1,2)\)，\(V_4''=(2,2)\).
由于 \(V_2''\) 与 \(V_4''\) 重合，俯视图的外轮廓是三角形 \(V_1''V_3''V_2''\)，与图\(\circled{2}\)一致. 因此正视图和俯视图分别为图\(\circled{4}\)和图\(\circled{2}\)，故选 D.
\end{solution}

\begin{problem}
若函数 \(f(x)\)，\(g(x)\) 满足 \( \int_{-1}^{1} f(x) g(x) \, \mathrm{d}x = 0 \)，则称 \(f(x)\)，\(g(x)\) 为区间 \([-1, 1]\) 上的一组正交函数，给出三组函数：

\begin{circlelist}
\item \(f(x)=\sin\frac{x}{2}\)，\(g(x)=\cos\frac{x}{2}\)；
\item \(f(x)=x+1\)，\(g(x)=x-1\)；
\item \(f(x)=x\)，\(g(x)=x^{2}\).
\end{circlelist}

其中为区间 \([-1, 1]\) 上的正交函数的组数是（\quad）

\choices
    {\(0\)}
    {\(1\)}
    {\(2\)}
    {\(3\)}
\end{problem}

\begin{answer}
C
\end{answer}

\begin{solution}
第 1 组函数的乘积为 \(\sin\frac{x}{2}\cos\frac{x}{2}=\frac12\sin x\)，故
\[
\int_{-1}^{1}f(x)g(x)\,\mathrm{d}x=\frac12\int_{-1}^{1}\sin x\,\mathrm{d}x=0
\]
第 2 组函数的乘积为 \(x^2-1\)，积分为 \(\left[\frac{x^3}{3}-x\right]_{-1}^{1}=-\frac43\ne0\). 第 3 组函数的乘积为 \(x^3\)，它是奇函数，在关于原点对称的区间上的积分为 \(0\). 所以共有 \(2\) 组，故选 C.
\end{solution}

\begin{problem}
由不等式组 \(\begin{cases}
x \leqslant 0,\\
y \geqslant 0,\\
y - x - 2 \leqslant 0
\end{cases}\) 确定的平面区域记为 \(\Omega_1\)，不等式组 \(\begin{cases}
x + y \leqslant 1,\\
x + y \geqslant -2
\end{cases}\) 确定的平面区域记为 \(\Omega_2\)，在 \(\Omega_1\) 中随机取一点，则该点恰好在 \(\Omega_2\) 内的概率为（\quad）

\choices
    {\(\frac{1}{8}\)}
    {\(\frac{1}{4}\)}
    {\(\frac{3}{4}\)}
    {\(\frac{7}{8}\)}
\end{problem}

\begin{answer}
D
\end{answer}

\begin{solution}
由 \(x\leqslant0\)、\(y\geqslant0\)、\(y\leqslant x+2\)，可知 \(-2\leqslant x\leqslant0\)，区域 \(\Omega_1\) 是顶点为 \((-2,0)\)、\((0,0)\)、\((0,2)\) 的三角形，面积为 \(2\). 在该三角形内，条件 \(x+y>1\) 截出的部分是顶点为 \((0,1)\)、\((0,2)\)、\((-\frac12,\frac32)\) 的三角形，面积为 \(\frac14\). 由于在 \(\Omega_1\) 内恒有 \(x+y\geqslant-2\)，所以 \(\Omega_1\cap\Omega_2\) 的面积为 \(2-\frac14=\frac74\). 因而所求概率为
\[
\frac{\frac74}{2}=\frac78
\]
故选 D.
\end{solution}

\begin{problem}
《算数书》竹简于上世纪八十年代在湖北省江陵县张家山出土，这是我国现存最早的有系统的数学典籍，其中记载有求“囷盖”的术：置如其周，令相乘也. 又以高乘之，三十六成一. 该术相当于给出了由圆锥的底面周长 \(L\) 与高 \(h\)，计算其体积 \(V\) 的近似公式 \(V \approx \frac{1}{36} L^2 h\). 它实际上是将圆锥体积公式中的圆周率 \(\pi\) 近似取为 \(3\). 那么，近似公式 \(V \approx \frac{2}{75} L^2 h\) 相当于将圆锥体积公式中的圆周率 \(\pi\) 近似取为（\quad）

\choices
    {\(\frac{22}{7}\)}
    {\(\frac{25}{8}\)}
    {\(\frac{157}{50}\)}
    {\(\frac{355}{113}\)}
\end{problem}

\begin{answer}
B
\end{answer}

\begin{solution}
由圆锥底面周长 \(L=2\pi r\)，得 \(r=\frac{L}{2\pi}\). 因此圆锥体积为
\[
V=\frac13\pi r^2h=\frac13\pi\left(\frac{L}{2\pi}\right)^2h=\frac{1}{12\pi}L^2h
\]
若题给近似式中的系数为 \(\frac{2}{75}\)，则 \(\frac{1}{12\pi}=\frac{2}{75}\)，从而 \(\pi=\frac{75}{24}=\frac{25}{8}\). 故选 B.
\end{solution}

\begin{problem}
已知 \(F_{1}\)，\(F_{2}\) 是椭圆和双曲线的公共焦点，\(P\) 是它们的一个公共点，且 \(\angle F_{1}PF_{2} = \frac{\pi}{3}\)，则椭圆和双曲线的离心率的倒数之和的最大值为（\quad）

\choices
    {\(\frac{4\sqrt{3}}{3}\)}
    {\(\frac{2\sqrt{3}}{3}\)}
    {\(3\)}
    {\(2\)}
\end{problem}

\begin{answer}
A
\end{answer}

\begin{solution}
设 \(PF_1=u\)，\(PF_2=v\)，两曲线的公共焦半距为 \(c\). 由椭圆和双曲线的定义，分别有 \(u+v=2a\) 且 \(|u-v|=2a'\).
其中 \(a\) 和 \(a'\) 分别为椭圆和双曲线的长半轴、实半轴. 令 \(s=u+v\)，\(d=|u-v|\)，则两条曲线离心率的倒数之和为
\[
\frac{1}{e_1}+\frac{1}{e_2}=\frac{a}{c}+\frac{a'}{c}=\frac{s+d}{2c}
\]
又因为 \(F_1F_2=2c\)，且 \(\angle F_1PF_2=\frac\pi3\)，由余弦定理得
\[
4c^2=u^2+v^2-uv=\frac{s^2+3d^2}{4}
\]
即 \(s^2+3d^2=16c^2\). 由柯西不等式，
\[
(s+d)^2\leqslant(s^2+3d^2)\left(1+\frac13\right)=\frac{64c^2}{3}
\]
所以
\[
\frac{s+d}{2c}\leqslant\frac{4\sqrt{3}}{3}
\]
当 \(s=3d\) 时取等号. 此时可取 \(u=2d\)，\(v=d\)，从而
\(c=\frac{\sqrt{3}}{2}d\)，\(a=\frac{3}{2}d\)，\(a'=\frac{1}{2}d\)，满足椭圆的 \(a>c\) 和双曲线的 \(c>a'\)，故等号可以取得. 因此最大值为 \(\frac{4\sqrt{3}}{3}\)，选 A.
\end{solution}

\begin{problem}
已知函数 \( f(x) \) 是定义在 \(\R\) 上的奇函数，当 \( x \geqslant 0 \) 时，\( f(x) = \frac{1}{2} (|x - a^2| + |x - 2a^2| - 3a^2) \). 若 \(\forall x \in \R\)，\(f(x - 1) \leqslant f(x)\)，则实数 \( a \) 的取值范围为（\quad）

\choices
    {\( \left[-\frac{1}{6},\frac{1}{6}\right] \)}
    {\( \left[-\frac{\sqrt{6}}{6},\frac{\sqrt{6}}{6}\right] \)}
    {\( \left[-\frac{1}{3},\frac{1}{3}\right] \)}
    {\( \left[-\frac{\sqrt{3}}{3},\frac{\sqrt{3}}{3}\right] \)}
\end{problem}

\begin{answer}
B
\end{answer}

\begin{solution}
令 \(t=a^2\geqslant0\)，并记当 \(u\geqslant0\) 时右侧给出的函数为 \(g(u)\)，则
\[
g(u)=\frac12\left(|u-t|+|u-2t|-3t\right)=
\begin{cases}
-u, & 0\leqslant u\leqslant t,\\
-t, & t\leqslant u\leqslant2t,\\
u-3t, & u\geqslant2t.
\end{cases}
\]
由奇性，\(f(-u)=-g(u)\). 当 \(0\leqslant x\leqslant1\) 时，题设不等式等价于 \(g(x)+g(1-x)\geqslant0\). 取 \(x=\frac12\)，若 \(t>\frac16\)，则分三种情形均有 \(g(\frac12)<0\)：当 \(t\leqslant\frac14\) 时，\(g(\frac12)=\frac12-3t<0\)；当 \(\frac14<t\leqslant\frac12\) 时，\(g(\frac12)=-t<0\)；当 \(t>\frac12\) 时，\(g(\frac12)=-\frac12<0\). 因此必须有 \(t\leqslant\frac16\).

反之，设 \(0\leqslant t\leqslant\frac16\). 对任意 \(s\geqslant0\)，由上式可直接看出 \(g(s+1)\geqslant g(s)\)：在 \([0,t]\) 上，差值为 \(2s+1-3t\geqslant0\)；在 \([t,2t]\) 上，差值为 \(s+1-2t\geqslant0\)；在 \([2t,+\infty)\) 上，\(g\) 单调递增. 这覆盖了 \(x\geqslant1\) 及 \(x\leqslant0\) 的情形.

再由对称性只需考察 \(0\leqslant x\leqslant\frac12\). 此时 \(1-x\geqslant\frac12\geqslant2t\)，故 \(g(1-x)=1-x-3t\). 当 \(0\leqslant x\leqslant t\) 时，\(g(x)+g(1-x)=1-2x-3t\geqslant1-5t\geqslant0\)；当 \(t\leqslant x\leqslant2t\) 时，该和为 \(1-x-4t\geqslant1-6t\geqslant0\)；当 \(2t\leqslant x\leqslant\frac12\) 时，该和为 \(1-6t\geqslant0\). 所以题设不等式成立当且仅当 \(t\leqslant\frac16\)，即
\[
|a|\leqslant\sqrt{\frac16}=\frac{\sqrt6}{6}
\]
故选 B.
\end{solution}

\section{填空题}

\begin{problem}
设向量 \(\bs{a} = (3, 3)\)，\(\bs{b} = (1, -1)\)，若 \((\bs{a} + \lambda \bs{b}) \perp (\bs{a} - \lambda \bs{b})\)，则实数 \(\lambda =\)\fillinblank{}.
\end{problem}

\begin{answer}
\(\lambda=\pm 3\)
\end{answer}

\begin{solution}
由垂直条件，
\[
(\bs{a}+\lambda\bs{b})\cdot(\bs{a}-\lambda\bs{b})
=|\bs{a}|^2-\lambda^2|\bs{b}|^2=0
\]
又 \(|\bs{a}|^2=3^2+3^2=18\)，\(|\bs{b}|^2=1^2+(-1)^2=2\)，所以 \(18-2\lambda^2=0\)，即 \(\lambda^2=9\). 因而 \(\lambda=\pm3\).
\end{solution}

\begin{problem}
直线 \(l_{1}: y = x + a\) 和 \(l_{2}: y = x + b\) 分别将单位圆 \(C: x^{2} + y^{2} = 1\) 分成长度相等的四段弧，则 \(a^{2} + b^{2}=\)\fillinblank{}.
\end{problem}

\begin{answer}
2
\end{answer}

\begin{solution}
每条直线的弦长所对应的圆心角必须为 \(\frac{\pi}{2}\)，因为四个交点把圆周分成四段等弧. 单位圆圆心到弦 \(y=x+a\) 的距离为 \(\frac{|a|}{\sqrt2}\)，而圆心到圆心角为 \(\frac{\pi}{2}\) 的弦的距离为
\[
\cos\left(\frac{\pi}{4}\right)=\frac{\sqrt2}{2}
\]
因此 \(\frac{|a|}{\sqrt2}=\frac{\sqrt2}{2}\)，得到 \(|a|=1\). 两条直线必须位于圆心两侧，否则它们不会给出四个依次相隔四分之一圆周的交点，所以 \(\{a,b\}=\{1,-1\}\). 从而
\[
a^2+b^2=1+1=2
\]
\end{solution}

\begin{problem}
设 \(a\) 是一个各位数字都不是 0 且没有重复数字的三位数. 将组成 \(a\) 的 \(3\text{ 个}\)数字按从小到大排成的三位数记为 \(I(a)\)，按从大到小排成的三位数记为 \(D(a)\)（例如 \(a = 815\)，则 \(I(a) = 158\)，\(D(a) = 851\)）. 阅读如图所示的程序框图，运行相应的程序，任意输入一个 \(a\)，输出的结果是 \(b =\)\fillinblank{}.

\texfigure[width=0.36\linewidth]{img_repaint/2014/hubei_science/q13_fig1.tex}
\end{problem}

\begin{answer}
495
\end{answer}

\begin{solution}
设输入三位数的三个数字从小到大为 \(p<q<r\). 第一次运算得到
\[
D(a)-I(a)=(100r+10q+p)-(100p+10q+r)=99(r-p)
\]
由于 \(p<q<r\) 且三个数字均不为 \(0\)，有 \(r-p\in\{2,3,\ldots,8\}\)，所以第一次得到的 \(b\) 属于
\[
\{198,297,396,495,594,693,792\}
\]
对这些可能值，程序中的连续结果分别为
\[
\begin{gathered}
198\to792\to693\to594\to495,\\
297\to693\to594\to495,\\
396\to594\to495,\\
495\to495,\\
594\to495,\\
693\to594\to495,\\
792\to693\to594\to495.
\end{gathered}
\]
又 \(D(495)-I(495)=954-459=495\)，所以程序最终必输出 \(b=495\).
\end{solution}

\begin{problem}
设 \( f(x) \) 是定义在 \((0, +\infty)\) 上的函数，且 \( f(x) > 0 \)，对任意 \(a > 0\)，\(b > 0\)，若经过点 \((a, f(a))\)，\((b, -f(b))\) 的直线与 \( x \) 轴的交点为 \((c, 0)\)，则称 \( c \) 为 \(a\)，\(b\) 关于函数 \( f(x) \) 的平均数，记为 \( M_{f}(a,b) \). 例如，当 \( f(x)=1 \)（\(x>0\)）时，可得 \( M_{f}(a,b)=c=\frac{a+b}{2} \)，即 \( M_{f}(a,b) \) 为 \(a\)，\(b\) 的算术平均数.
\begin{enumerate}
\item 当 \( f(x)=\)\fillinblank{}（\(x>0\)）时，\( M_f(a, b) \) 为 \(a\)，\(b\) 的几何平均数；
\item 当 \( f(x)=\)\fillinblank{}（\(x>0\)）时，\( M_f(a, b) \) 为 \(a\)，\(b\) 的调和平均数 \(\frac{2ab}{a + b}\). （以上两空各只需写出一个符合要求的函数即可）
\end{enumerate}
\end{problem}

\begin{answer}
\(f(x)=\sqrt{x}\)；\(f(x)=x\)
\end{answer}

\begin{solution}
设直线经过 \((a,f(a))\) 和 \((b,-f(b))\)，其与 \(x\) 轴交于 \((c,0)\). 由截距公式，
\[
c=\frac{a f(b)+b f(a)}{f(a)+f(b)}
\]
当 \(f(x)=\sqrt{x}\) 时，
\[
c=\frac{a\sqrt b+b\sqrt a}{\sqrt a+\sqrt b}=\sqrt{ab}
\]
所以得到几何平均数. 当 \(f(x)=x\) 时，
\[
c=\frac{ab+ba}{a+b}=\frac{2ab}{a+b}
\]
所以得到调和平均数. 这两个函数均满足题设的正值条件.
\end{solution}

\section{选考题：填空题}

\begin{problem}
如图，\(P\) 为 \( \odot O \) 外一点，过 \(P\) 点作 \( \odot O \) 的两条切线，切点分别为 \(A\)，\(B\)，过 \(PA\) 的中点 \(Q\) 作割线交 \( \odot O \) 于 \(C\)，\(D\) 两点. 若 \(QC = 1\)，\(CD = 3\)，则 \(PB =\)\fillinblank{}.

\bitmapfigure[max width=0.42\linewidth]{img/2014/hubei_science/q15_fig1.png}
\end{problem}

\begin{answer}
4
\end{answer}

\begin{solution}
因为 \(PA\) 是圆的切线，且 \(Q\) 在 \(PA\) 上，所以 \(QA\) 也是从 \(Q\) 引圆的切线. 由点 \(Q\) 的切割线定理，
\[
QA^2=QC\cdot QD=1\cdot(1+3)=4
\]
从而 \(QA=2\). 又 \(Q\) 是 \(PA\) 的中点，故 \(PA=2QA=4\). 两切线段相等，所以 \(PB=PA=4\).
\end{solution}

\begin{problem}
已知曲线 \(C_{1}\) 的参数方程是 \(\begin{cases}
x = \sqrt{t},\\
y = \frac{\sqrt{3t}}{3},
\end{cases}\)（\(t\) 为参数），以坐标原点为极点，\(x\) 轴的正半轴为极轴建立极坐标系，曲线 \(C_{2}\) 的极坐标方程是 \(\rho = 2\)，则 \(C_{1}\) 与 \(C_{2}\) 交点的直角坐标为\fillinblank{}.
\end{problem}

\begin{answer}
\((\sqrt{3},1)\)
\end{answer}

\begin{solution}
由 \(x=\sqrt t\) 可知 \(t=x^2\)，且 \(x\geqslant0\). 因而
\[
y=\frac{\sqrt{3t}}{3}=\frac{x}{\sqrt3}
\]
即曲线 \(C_1\) 的直角坐标方程为 \(y=\frac{x}{\sqrt3}\)，\(x\geqslant0\). 极坐标方程 \(\rho=2\) 对应的直角坐标方程为 \(x^2+y^2=4\). 联立得
\[
x^2+\frac{x^2}{3}=4
\]
所以 \(x^2=3\). 由 \(x\geqslant0\)，得 \(x=\sqrt3\)，再得 \(y=1\). 故交点坐标为 \((\sqrt3,1)\).
\end{solution}

\section{解答题}

\begin{problem}
某实验室一天的温度（单位：\(^{\circ}\mathrm{C}\)）随时间 \(t\)（单位：\(\mathrm{h}\)）的变化近似满足函数关系： \(f(t) = 10 - \sqrt{3}\cos \frac{\pi}{12}t - \sin \frac{\pi}{12}t\)，\(\ t \in [0, 24)\).
\begin{enumerate}
\item 求实验室这一天的最大温差；
\item 若要求实验室温度不高于 \(11 \mathrm{~^{\circ}C}\)，则在哪段时间实验室需要降温？
\end{enumerate}
\end{problem}

\begin{solution}
\begin{enumerate}
\item 令 \(u=\frac{\pi}{12}t\)，则 \(0\leqslant u<2\pi\)，且
\[
\sqrt3\cos u+\sin u=2\cos\left(u-\frac{\pi}{6}\right)
\]
所以 \(f(t)=10-2\cos\left(u-\frac{\pi}{6}\right)\)，从而 \(8\leqslant f(t)\leqslant12\). 这一天的最大温差为 \(12-8=4 \mathrm{~^{\circ}C}\).
\item 当温度需要降到不高于 \(11 \mathrm{~^{\circ}C}\) 时，需在 \(f(t)>11\) 的时段降温. 由
\[
10-2\cos\left(u-\frac{\pi}{6}\right)>11
\]
得 \(\cos\left(u-\frac{\pi}{6}\right)<-\frac12\). 在 \(0\leqslant u<2\pi\) 内，这等价于 \(\frac{5\pi}{6}<u<\frac{3\pi}{2}\). 因为 \(u=\frac{\pi}{12}t\)，所以 \(10<t<18\). 即实验室在第\(10\text{ 小时}\)至第\(18\text{ 小时}\)之间需要降温.
\end{enumerate}
\end{solution}

\begin{problem}
已知等差数列 \(\{a_{n}\}\) 满足： \(a_{1} = 2\)，且 \(a_{1}\)，\(a_{2}\)，\(a_{5}\) 成等比数列.
\begin{enumerate}
\item 求数列 \(\{a_{n}\}\) 的通项公式.
\item 记 \(S_{n}\) 为数列 \(\{a_{n}\}\) 的前 \(n\) 项和，是否存在正整数 \(n\)，使得 \(S_{n} > 60n + 800\)？若存在，求出 \(n\) 的最小值；若不存在，说明理由.
\end{enumerate}
\end{problem}

\begin{solution}
设等差数列的公差为 \(d\)，则 \(a_2=2+d\)，\(a_5=2+4d\). 由 \(a_1\)，\(a_2\)，\(a_5\) 成等比数列，有
\[
(2+d)^2=2(2+4d)
\]
化简得 \(d^2-4d=0\)，所以 \(d=0\) 或 \(d=4\).

\begin{enumerate}
\item 当 \(d=0\) 时，\(a_n=2\)；当 \(d=4\) 时，\(a_n=2+4(n-1)=4n-2\). 因而数列的通项公式有两种可能：
\[
a_n=2\quad\text{或}\quad a_n=4n-2
\]
\item 当 \(d=0\) 时，\(S_n=2n<60n+800\)，不存在满足条件的正整数 \(n\). 当 \(d=4\) 时，
\[
S_n=\frac{n(a_1+a_n)}{2}=\frac{n(2+4n-2)}{2}=2n^2
\]
于是 \(2n^2>60n+800\)，即 \(n^2-30n-400>0\)，也就是 \((n-40)(n+10)>0\). 因为 \(n\) 是正整数，故 \(n\geqslant41\)，最小值为 \(41\). 综上，\(a_n=2\) 时不存在；\(a_n=4n-2\) 时存在，最小的 \(n\) 为 \(41\).
\end{enumerate}
\end{solution}

\begin{problem}
如图，在棱长为 2 的正方体 \(ABCD - A_{1}B_{1}C_{1}D_{1}\) 中，\(E\)，\(F\)，\(M\)，\(N\) 分别是棱 \(AB\)，\(AD\)，\(A_{1}B_{1}\)，\(A_{1}D_{1}\) 的中点，点 \(P\)，\(Q\) 分别在棱 \(DD_{1}\)，\(BB_{1}\) 上移动，且 \(DP = BQ = \lambda\)（\(0 < \lambda < 2\)）.
\begin{enumerate}
\item 当 \(\lambda = 1\) 时，证明：直线 \(BC_{1}\) \(\parallel\) 平面 \(EFPQ\).
\item 是否存在 \(\lambda\)，使平面 \(EFPQ\) 与面 \(PQMN\) 所成的二面角为直二面角？若存在，求出 \(\lambda\) 的值；若不存在，说明理由.
\end{enumerate}

\bitmapfigure[max width=0.62\linewidth]{img/2014/hubei_science/q19_fig1.png}
\end{problem}

\begin{solution}
建立空间直角坐标系，取 \(A(0,0,0)\)，\(B(2,0,0)\)，\(D(0,2,0)\)，\(A_1(0,0,2)\). 则 \(E(1,0,0)\)，\(F(0,1,0)\)，\(M(1,0,2)\)，\(N(0,1,2)\)，\(P(0,2,\lambda)\)，\(Q(2,0,\lambda)\)，\(C_1(2,2,2)\).
\begin{enumerate}
\item 当 \(\lambda=1\) 时，\(\overrightarrow{EF}=(-1,1,0)\)，\(\overrightarrow{EP}=(-1,2,1)\)，且
\[
\overrightarrow{BC_1}=(0,2,2)=2\left(\overrightarrow{EP}-\overrightarrow{EF}\right)
\]
因此直线 \(BC_1\) 的方向向量是平面 \(EFPQ\) 内两个相交直线 \(EF\)、\(EP\) 的方向向量的线性组合，故 \(BC_1\) 的方向平行于平面 \(EFPQ\). 另外，\(\overrightarrow{EF}\times\overrightarrow{EP}=(1,1,-1)\)，所以平面 \(EFPQ\) 的方程为 \(x+y-z=1\)，而 \(B(2,0,0)\) 不在此平面上. 因此 \(BC_1\parallel\) 平面 \(EFPQ\).
\item 取平面 \(EFPQ\) 的两个方向向量 \(\overrightarrow{EF}=(-1,1,0)\)、\(\overrightarrow{EP}=(-1,2,\lambda)\)，可得其一个法向量为
\[
\bs{n}_1=(\lambda,\lambda,-1)
\]
平面 \(PQMN\) 内可取 \(\overrightarrow{PQ}=(2,-2,0)\)、\(\overrightarrow{PN}=(0,-1,2-\lambda)\)，故其一个法向量可取
\[
\bs{n}_2=(\lambda-2,\lambda-2,-1)
\]
两个平面所成的二面角为直二面角，当且仅当两个法向量垂直，即
\[
\bs{n}_1\cdot\bs{n}_2=2\lambda(\lambda-2)+1=0
\]
解得
\[
\lambda=1\pm\frac{\sqrt2}{2}
\]
这两个值都属于 \((0,2)\)，所以存在，且所求值为 \(\lambda=1-\frac{\sqrt2}{2}\) 或 \(\lambda=1+\frac{\sqrt2}{2}\).
\end{enumerate}
\end{solution}

\begin{problem}
计划在某水库建一座至多安装 \(3\text{ 台}\)发电机的水电站，过去 \(50\text{ 年}\)的水文资料显示，水库年入流量 \(X\)（年入流量：一年内上游来水与库区降水之和，单位：\(\text{亿立方米}\)）都在 \(40\) 以上. 其中，不足 \(80\) 的年份有 \(10\text{ 年}\)，不低于 \(80\) 且不超过 \(120\) 的年份有 \(35\text{ 年}\)，超过 \(120\) 的年份有 \(5\text{ 年}\). 将年入流量在以上三段的频率作为相应段的概率，并假设各年的年入流量相互独立.
\begin{enumerate}
\item 求未来 \(4\text{ 年}\)中，至多有 \(1\text{ 年}\)的年入流量超过 \(120\) 的概率；
\item 水电站希望安装的发电机尽可能运行，但每年发电机最多可运行台数受年入流量 \(X\) 限制，并有如下关系：

\begin{center}
\begin{tabular}{c|ccc}
\hline
年入流量 \(X\) & \(40<X<80\) & \(80\leqslant X\leqslant 120\) & \(X>120\) \\
\hline
发电机最多可运行台数 & \(1\) & \(2\) & \(3\) \\
\hline
\end{tabular}
\end{center}

若某台发电机运行，则该台年利润为 \(5000\text{ 万元}\)；若某台发电机未运行，则该台年亏损 \(800\text{ 万元}\)，欲使水电站年总利润的均值达到最大，应安装发电机多少台？
\end{enumerate}
\end{problem}

\begin{solution}
\begin{enumerate}
\item 设一年中年入流量超过 \(120\) 的事件为 \(A\)，则 \(P(A)=\frac{5}{50}=\frac{1}{10}\). 未来四年相互独立，所以至多一年发生事件 \(A\) 的概率为
\[
\left(\frac9{10}\right)^4+\mathrm C_4^1\frac1{10}\left(\frac9{10}\right)^3
=\frac{9^4+4\cdot9^3}{10^4}
=\frac{9477}{10000}
\]
\item 安装 \(m\) 台发电机时，每年实际运行台数为 \(\min(m,K)\)，其中 \(K\) 取 \(1\)，\(2\)，\(3\) 的概率分别为 \(\frac15\)，\(\frac7{10}\)，\(\frac1{10}\). 逐一计算三种安装方案的年总利润均值：
\[
E_1=5000
\]
\[
E_2=\frac15(5000-800)+\frac7{10}(2\cdot5000)+\frac1{10}(2\cdot5000)=8840
\]
\[
E_3=\frac15(5000-2\cdot800)+\frac7{10}(2\cdot5000-800)+\frac1{10}(3\cdot5000)=8620
\]
其中 \(E_2\) 最大，所以应安装 \(2\) 台发电机.
\end{enumerate}
\end{solution}

\begin{problem}
在平面直角坐标系 \(xOy\) 中，点 \(M\) 到点 \(F(1,0)\) 的距离比它到 \(y\) 轴的距离多 1，记点 \(M\) 的轨迹为 \(C\).
\begin{enumerate}
\item 求轨迹为 \(C\) 的方程；
\item 设斜率为 \(k\) 的直线 \(l\) 过定点 \(P(-2,1)\)，求直线 \(l\) 与轨迹 \(C\) 恰好有一个公共点，两个公共点，三个公共点时 \(k\) 的相应取值范围.
\end{enumerate}
\end{problem}

\begin{solution}
\begin{enumerate}
\item 设 \(M(x,y)\). 由题意，
\[
\sqrt{(x-1)^2+y^2}=|x|+1
\]
两边平方并化简得 \(y^2=2x+2|x|\). 当 \(x\geqslant0\) 时，\(y^2=4x\)；当 \(x<0\) 时，\(y=0\). 反代可知这些点均满足原距离关系，因此
\(C:\)，\(y^2=4x\ (x\geqslant0)\quad\text{或}\quad y=0\ (x\leqslant0)\).
\item 直线 \(l\) 的方程为 \(y=kx+2k+1\). 轨迹的射线部分 \(y=0\)，\(x\leqslant0\) 与直线的交点在 \(k\ne0\) 时满足 \(x=-2-\frac1k\)，所以当 \(k\leqslant-\frac12\) 或 \(k>0\) 时有一个交点，当 \(-\frac12<k\leqslant0\) 时没有交点.

对于抛物线部分，令 \(x=\frac{y^2}{4}\)，代入直线方程得
\[
ky^2-4y+4(2k+1)=0
\]
当 \(k\ne0\) 时，其判别式为 \(16(1-k-2k^2)\). 因此在 \(k\in(-1,\frac12)\) 且 \(k\ne0\) 时有两个交点，在 \(k=-1\) 或 \(k=\frac12\) 时有一个交点，在区间外没有交点. 当 \(k=0\) 时方程退化为 \(y=1\)，有一个交点.

特别地，\(k=-\frac12\) 时直线经过原点，原点同时属于射线和抛物线，计数时只能算一个公共点. 综合上述各情形，恰有一个、两个、三个公共点时分别有
\(k\in(-\infty,-1)\cup\{0\}\cup\left(\frac12,+\infty\right)\) 时恰有一个公共点；\(k\in\{-1\}\cup\left[-\frac12,0\right)\cup\left\{\frac12\right\}\) 时恰有两个公共点；\(k\in\left(-1,-\frac12\right)\cup\left(0,\frac12\right)\) 时恰有三个公共点.
\end{enumerate}
\end{solution}

\begin{problem}
\(\pi\) 为圆周率，\(\e = 2.71828\cdots\) 为自然对数的底数.
\begin{enumerate}
\item 求函数 \( f(x) = \frac{\ln x}{x} \) 的单调区间；
\item 求 \( \e^{3} \)，\( 3^{\e} \)，\( \e^{\pi} \)，\( \pi^{\e} \)，\( 3^{\pi} \)，\( \pi^{3} \) 这 6 个数中的最大数与最小数；
\item 将 \( \e^{3} \)，\( 3^{\e} \)，\( \e^{\pi} \)，\( \pi^{\e} \)，\( 3^{\pi} \)，\( \pi^{3} \) 这 6 个数按从小到大的顺序排列，并证明你的结论.
\end{enumerate}
\end{problem}

\begin{solution}
\begin{enumerate}
\item
\[
f'(x)=\frac{1-\ln x}{x^2}
\]
因为 \(x^2>0\)，所以当 \(0<x<\e\) 时 \(f'(x)>0\)，当 \(x>\e\) 时 \(f'(x)<0\). 因此 \(f(x)\) 在 \((0,\e)\) 上单调递增，在 \((\e,+\infty)\) 上单调递减.
\item 下面比较这些数的自然对数. 由 \(f(x)=\frac{\ln x}{x}\) 在 \((\e,+\infty)\) 上递减，且 \(\e<3<\pi\)，有
 \(\frac{\ln 3}{3}<\frac{1}{\e}\)，\(\frac{\ln \pi}{\pi}<\frac{1}{\e}\)，\(\frac{\ln \pi}{\pi}<\frac{\ln 3}{3}\).
这分别给出 \(3^{\e}<\e^{3}\)、\(\pi^{\e}<\e^{\pi}\) 和 \(\pi^{3}<3^{\pi}\). 为比较中间两对，对 \(x>0\) 令 \(z=\frac{x-1}{x+1}\)，则 \(-1<z<1\)，且有展开式
\[
\ln x=2\left(z+\frac{z^3}{3}+\frac{z^5}{5}+\cdots\right)
\]
当 \(x>1\) 时 \(z>0\)，所以截取前三项可得下界.
由常用估值 \(3.14<\pi<\frac{22}{7}\) 及 \(\e>2.718\)，有
\[
\ln\pi>\ln 3.14>2\left(\frac{107}{207}+\frac{1}{3}\left(\frac{107}{207}\right)^3+\frac{1}{5}\left(\frac{107}{207}\right)^5\right)>1.14
\]
\[
\frac{3}{\e}<\frac{3}{2.718}<1.104
\]
所以 \(\frac{3}{\e}<\ln\pi\)，即 \(3<\e\ln\pi\)，从而 \(\e^{3}<\pi^{\e}\). 另一方面，利用 \(\pi>3\)，有
\[
\ln\pi>\ln3>2\left(\frac12+\frac{1}{3\cdot2^3}+\frac{1}{5\cdot2^5}\right)=\frac{263}{240}>\frac{22}{21}>\frac{\pi}{3}
\]
故 \(3\ln\pi>\pi\)，从而 \(\e^{\pi}<\pi^{3}\).
因此这 6 个数中最小数为 \(3\) 的 \(\e\) 次方，最大数为 \(3\) 的 \(\pi\) 次方.

\item 将上面的严格不等式依次连接，得到如下从小到大的排列：
\[
3^{\e}<\e^{3}<\pi^{\e}<\e^{\pi}<\pi^{3}<3^{\pi}
\]
\end{enumerate}
\end{solution}
